% Cote 119: offprints, reports, courses, CV and letters (typed and printed part): archive cote 119, site page 226 (PDF page 227).
% typescript. Draft compared: none.
% AI-transcribed by Claude Opus 5.5 (claude-opus-5-5, adaptive thinking, effort high), Message Batches API, batch r119_full (msgbatch_01DwZNm5Y7YsKBWyAN4sFs8w), 2026-11-13; not yet checked by eye.
\scanpage{119/226}
\mspage{. 13 .}

et leur importance a été soulignée à de nombreuses
occasions par des mathématiciens aussi divers que
O. Zariski, J.P. Serre, H. Hironaka, D. Mumford, I. Manin,
F. Chafarévitch.

b) \uline{Théorie locale des schémas et des morphismes de schémas} :
Dans ce contexte se placent les développements d'algèbre
commutative mentionnés au par. 4, et l'étude détaillée
de notions comme celles de morphisme lisse, étale,
net, plat, etc. Les quatre volumes de [15, Chap. IV]
sont consacrés à ces développements, qui ont d'ailleurs
inspiré des développements analogues en théorie des
espaces analytiques et rigide-analytiques.

c) \uline{Techniques de construction de schémas} : Parmi les
techniques développées, exposées surtout dans [15 bis]
et des séminaires non publiés (par moi-même et d'autres),
il y a la \uline{théorie de la descente} (cf. aussi [16, SGA I,
Exp. V, VI]), celle des \uline{schémas quotients}, des \uline{schémas}
\uline{de Hilbert}, des \uline{schémas de Picard}, des \uline{``modules'' formels},
le \uline{théorème d'existence} des faisceaux de modules algébri-%
ques associés à des modules formels ([15, Chap. III,
par. 5]). Le point de vue adopté est surtout celui de
la construction d'un schéma à partir du foncteur qu'il
représente. Dans cette optique, je n'étais pas parvenu
à une véritable caractérisation maniable des foncteurs
représentables par un schéma relatif (localement de type
fini sur un schéma noethérien) - c'est M. Artin qui y
est parvenu ultérieurement [$*$], en remplaçant la notion
de schéma par celle, plus générale et plus stable,
d'espace algébrique. Parmi d'autres recherches dans
la même direction, suscitées par mes travaux, il y a
celles de P. Murre sur les schémas de Picard sur un
corps [$*$], celles de D. Mumford et de M. Raynaud [$*$]

\margin{226}
